\subsubsection{PHÉP ĐO CÁC ĐẠI LƯỢNG VẬT LÍ – HỆ ĐƠN VỊ SI}
\paragraph{Định nghĩa phép đo một đại lượng vật lí}
\newghichu{\textit{Phép đo một đại lượng vật lí là phép so sánh nó với một đại lượng cùng loại được quy ước làm đơn vị.}} 
\paragraph{Đơn vị đo}
\begin{itemize}
\item Hệ đơn vị đo thông dụng hiện nay là hệ SI.
\item Hệ SI quy định \textbf{7 đơn vị cơ bản}:
\end{itemize}
\resizebox{\textwidth}{!}{%
\begin{tabular}{|c|c|c|c|c|c|c|}
\hline 
Độ dài & Thời gian & Khối lượng & Nhiệt độ & Cường độ dòng điện & Cường độ sáng & Lượng chất \\ 
\hline 
mét (m) & giây (s) & kilôgam (kg) & kelvin (K) & ampe (A) & candela (Cd) & mol (mol) \\ 
\hline 
\end{tabular}%
}
\begin{itemize}
\item Ngoài 7 đơn vị cơ bản, các đơn vị còn lại được gọi là \textbf{đơn vị dẫn xuất}.
\end{itemize}
%\paragraph{Thứ nguyên}
%\newghichu{\textit{Tên gọi chung của tất cả các đơn vị trong các hệ đo lường khác nhau của cùng một đại lượng được gọi là \textbf{thứ nguyên} của đại lượng đó}}
%\begin{itemize}
%	\item Mỗi đại lượng vật lí chỉ có một thứ nguyên duy nhất.
%	\item Một số đại lượng vật lí có thể có cùng thứ nguyên.
%	\item Thứ nguyên của một đại lượng X được biểu diễn dưới dạng [X].
%	\item Thứ nguyên của một số đại lượng cơ bản:
%\end{itemize}
%\begin{center}
%	$\begin{array}{|c|c|}
%	\hline \begin{array}{c}
%		\text { \textbf{Đại lượng cơ bản}} 
%	\end{array} & \begin{array}{c}
%		\text { \textbf{Thứ nguyên} }
%	\end{array} \\
%	\hline \text { [Chiều dài] } & L \\
%	\hline \text { [Khối lượng] } & M \\
%	\hline \text { [Thời gian] } & T \\
%	\hline \text { [Cường độ dòng điện] } & I \\
%	\hline \text { [Nhiệt độ] } & K \\
%	\hline
%\end{array}$
%\end{center}
\paragraph{Kí hiệu bội số và ước số của đơn vị đo}
\begin{itemize}
\item Khi số đo của đại lượng là một bội số hoặc ước số thập phân của 10, ta có thể sử dụng kí hiệu (\textit{tiếp đầu ngữ}) ngay trước đơn vị để phần số đo được trình bày ngắn gọn.
\item Tên và kí hiệu một số tiếp đầu ngữ:
\end{itemize}
\begin{center}
$ \begin{array}{|c|c|c|c|c|c|}
\hline \begin{array}{c}
\text { Kí hiệu }
\end{array} & \begin{array}{l}
\text { Tên đọc }
\end{array} & \begin{array}{l}
\text { Hệ số }
\end{array} & \begin{array}{c}
\text { Ki hiệu }
\end{array} & \begin{array}{l}
\text { Tên đọc }
\end{array} & \begin{array}{l}
\text { Hệ số }
\end{array} \\
\hline \mathrm{T} & \text { tera } & 10^{12} & \text { d } & \text { deci } & 10^{-1} \\
\hline \text { G } & \text { giga } & 10^9 & \text { c } & \text { centi } & 10^{-2} \\
\hline \text { M } & \text { mega } & 10^6 & \mathrm{~m} & \text { mili } & 10^{-3} \\
\hline \text { k } & \text { kilo } & 10^3 & \mu & \text { micro } & 10^{-6} \\
\hline \text { h } & \text { hecto } & 10^2 & \mathrm{n} & \text { nano } & 10^{-9} \\
\hline \mathrm{da} & \text { deka } & 10^1 & \mathrm{p} & \text { pico } & 10^{-12} \\
\hline
\end{array} $
\end{center}
\paragraph{Các phương pháp đo một đại lượng vật lí}
\begin{itemize}
\item Đo trực tiếp một đại lượng bằng dụng cụ đo, kết quả được đọc trực tiếp trên dụng cụ đo được gọi là \textit{phép đo trực tiếp}.
\item Đo một đại lượng không trực tiếp mà thông qua công thức liên hệ với các đại lượng có thể đo trực tiếp gọi là \textit{phép đo gián tiếp}.
\end{itemize}

\paragraph{Số chữ số có nghĩa trong kết quả đo}
Số chữ số có nghĩa là tất cả những chữ số tính từ trái sang phải kể từ chữ số khác 0 đầu tiên.
\begin{itemize}
\item $13,1$ $\xrightarrow{[13,1]}$ có 3 chữ số có nghĩa.
\item $13,10$ $\xrightarrow{[13,10]}$ có 4 chữ số có nghĩa.
\item $1,30.10^3$ $\xrightarrow{[1,30]}$ có 3 chữ số có nghĩa.
\item $0,3109$ $\xrightarrow{0,[3109]}$ có 4 chữ số có nghĩa.
\item $0,0314$ $\xrightarrow{0,0[314]}$ có 3 chữ số có nghĩa.
\end{itemize}
\subsubsection{SAI SỐ PHÉP ĐO}
\setcounter{paragraph}{0}
\paragraph{Phân loại sai số}
\begin{itemize}
\item[a)] \textbf{Sai số hệ thống}
\begin{itemize}
\item Khi sử dụng dụng cụ đo để đo các đại lượng vật lí luôn có sự sai lệch do đặc điểm và cấu tạo của dụng cụ gây ra. Sự sai lệch này gọi là sai số dụng cụ hoặc sai số hệ thống.
\item Sai số hệ thống có nguyên nhân khách quan (do dụng cụ), nguyên nhân chủ quan do người đo (cần loại bỏ).
\end{itemize}	
\item[b)] \textbf{Sai số ngẫu nhiên}
\begin{itemize}
\item Khi lặp lại các phép đo, ta nhận được các giá trị khác nhau, sự sai lệch này không có nguyên nhân rõ ràng nên gọi là sai số ngẫu nhiên, có thể do thao tác không chuẩn, điều kiện làm thí nghiệm không ổn định hoặc hạn chế về giác quan,...
\item Để khắc phục người ta thường tiến hành thí nghiệm nhiều lần và tính sai số.
\end{itemize}	
\end{itemize}
\begin{note}
Sai số gây bởi dụng cụ sử dụng đo có thể lấy bằng một độ chia nhỏ nhất hoặc nửa độ chia nhỏ nhất trên dụng cụ (ví dụ thước đo chiều dài có độ chia nhỏ nhất là 1 mm thì sai số dụng cụ là 0,5 mm), hoặc được ghi trực tiếp trên dụng cụ do nhà sản xuất xác định
\end{note}
\paragraph{Cách xác định sai số phép đo trực tiếp}
\begin{itemize}
\item \textit{\textbf{Sai số tuyệt đối của mỗi phép đo}} là trị tuyệt đối của hiệu số giữa giá trị trung bình các lần đo và giá trị của mỗi lần đo của phép đo trực tiếp.
\begin{align*}
\Delta A_1=\left|\overline{A}-{A}_1\right|;\ \Delta A_2=\left|\overline{A}-{A}_2\right|;...;\Delta A_n=\left|\overline{A}-{A}_n\right|
\end{align*}
Trong đó: $\overline{A}=\dfrac{A_1+A_2+...+A_n}{n}$\\
\textit{\textbf{Sai số tuyệt đối trung bình của n lần đo}} được tính theo công thức:
\begin{align*}
\overline{\Delta A}=\dfrac{\Delta A_1+\Delta A_2+...+\Delta A_n}{n}
\end{align*}
\textit{\textbf{Sai số tuyệt đối của phép đo}} là tổng sai số dụng cụ và sai số ngẫu nhiên:
\begin{align*}
\Delta A=\overline{\Delta A}+\Delta A_{dc}
\end{align*}
\item \textit{\textbf{Sai số tỉ đối của phép đo}} là tỉ lệ phần trăm giữa sai số tuyệt đối và giá trị trung bình của đại lượng đo, cho biết mức độ chính xác của phép đo.
\begin{align*}
\delta A=\dfrac{\Delta A}{\overline{A}}.100\%
\end{align*}
\end{itemize}
\paragraph{Cách xác định sai số phép đo gián tiếp}
\begin{itemize}
\item Phép đo trực tiếp: $X=\overline{X}\pm \Delta X;\quad \quad Y=\overline{Y}\pm \Delta Y;\quad \quad Z=\overline{Z}\pm \Delta Z$
\item Để xác định sai số của phép đo gián tiếp, vận dụng quy tắc sau:
\begin{itemize}
\item[$\bullet$] \textit{Sai số tuyệt đối của một tổng hay hiệu bằng tổng các sai số tuyệt đối của các số hạng}\\
Ví dụ ta có công thức: $A=2X + 3Y - 4Z$
$\Rightarrow$ 
$\begin{cases}
 \overline{A}=2\overline{X} + 3\overline{Y} - 4\overline{Z}\\ 
\Delta A=2\Delta X+3\Delta Y+4\Delta Z
\end{cases}$
\item[$\bullet$] \textit{Sai số tỉ đối của một tích hay thương thì bằng tổng các sai số tỉ đối của các thừa số}
\end{itemize}
\item Ví dụ ta có công thức: $A=3.\dfrac{X.Y^\alpha}{Z^\beta}$
$\Rightarrow$ 
$\begin{cases}
\overline{A}=3.\dfrac{\overline{X}.\overline{Y}^\alpha}{\overline{Z}^\beta}\\ 
\delta A=\delta X+\alpha.\delta Y+\beta.\delta Z
\end{cases}$
\end{itemize}
%\begin{note}
%	\begin{itemize}
%		\item Đo quãng đường s từ A đến C bằng tổng quãng đường $s_1$ từ A đến B và $s_2$ từ B đến C.\\
%		Sai số tuyệt đối: $\Delta s=\Delta s_1+\Delta s_2$
%		\item Đo tốc độ theo công thức $v=\dfrac{s}{t}$, sai số phép đo là:
%		\begin{align*}
%			&\delta v=\dfrac{\Delta s}{\overline{s}}.100\%+\dfrac{\Delta t}{\overline{t}}.100\%\\
%			&\Delta v=\delta v.\overline{v}
%		\end{align*}
%	\end{itemize}
%\end{note}
\paragraph{Cách ghi kết quả đo}
\begin{itemize}
\item Kết quả đo đại lượng A được ghi dưới dạng một khoảng giá trị:
\begin{align*}
\left(\overline{A}-\Delta A\right)\leq A \leq \left(\overline{A}+\Delta A\right)\ \text{hoặc}\ A=\overline{A}\pm \Delta A
\end{align*}
Trong đó:
\begin{itemize}
\item[+] $\Delta A$ là sai số tuyệt đối thường viết đến tối đa 2 chữ số có nghĩa.
\item[+] Phần thêm hoặc bớt không được vượt quá $ \dfrac{1}{10} $ giá trị $\overline{A}$.
\item[+] Giá trị trung bình $\overline{A}$ được viết đến bậc thập phân tương ứng với $\Delta A$.
\item[+] Nếu giá trị trung bình viết dưới dạng $ A.10^n $ thì sai số tuyệt đối cũng phải viết dạng $ B.10^n $.
\end{itemize}
\item Quy tắc làm tròn số
\begin{itemize}
\item[+] Nếu chữ số ở hàng bỏ đi nhỏ hơn 5 thì làm tròn xuống.
\item[+] Nếu chữ số ở hàng bỏ đi lớn hơn hoặc bằng 5 thì làm tròn lên. 
\end{itemize}
\end{itemize}
\begin{note}
\immini{\par Sai số và kết quả của phép đo có thể biểu diễn bằng đồ thị. Phương pháp đồ thị cho phép tìm quy luật sự phụ thuộc của đại lượng y vào đại lượng x.
\par Mỗi giá trị có được từ thực nghiệm đều có sai số: $\pm \Delta x$ và $\pm \Delta y$. Do đó trên đồ thị, mỗi giá trị sẽ được biểu diễn bằng một điểm nằm trong ô chữ nhật có cạnh là $2\Delta x$ và $2\Delta y$}{\begin{tikzpicture}[scale=1,thick,>=stealth',draw=Mapcolor]
\tkzInit[ymin=-0.5,ymax=4.5,xmin=-1,xmax=6]\tkzClip
\tkzDefPoints{0/0/o, 0/4/y, 5.5/0/x, 0/1/a}
\tkzDefShiftPoint[a](20:3cm){b}
\tkzDefShiftPoint[b](20:2.5cm){c}
\tkzDefPointBy[projection=onto o--x](b)
\tkzGetPoint{bx}
\tkzDefPointBy[projection=onto o--y](b) \tkzGetPoint{by}
\tkzDrawSegments[very thick,->](o,x o,y)
\tkzDrawSegments[very thick,Mapcolor](a,c)
\tkzDrawSegments[dashed,Mapcolor](b,bx b,by)
\tkzLabelPoint[below](o){$O$}
\tkzLabelPoint[below](x){$x$}
\tkzLabelPoint[left](y){$y$}
\tkzLabelPoint[below](bx){$x_M$}
\tkzLabelPoint[left](by){$y_M$}
\node  at  (b)  [shape=rectangle,draw=red,minimum height=0.9cm,minimum width=0.6cm, rounded corners=1pt,label=above:$2\Delta x$,label=right:$2\Delta y$]  {};
\node  at  (4.5,2.5)  [shape=rectangle,draw=red,minimum height=0.9cm,minimum width=0.6cm, rounded corners=1pt,label=above:Ô bao sai số]  {};
\node  at  (1,1.5)  [shape=rectangle,draw=red,minimum height=0.9cm,minimum width=0.6cm, rounded corners=1pt]  {};
\end{tikzpicture}}

\end{note}	






